
\Titular*%
{divulgacion}%
{Sobre a lóxica da física}%
{Ana Peón Nieto}%
{}%

\begin{multicols}{2}

A ninguén sorprenderá se digo que este é un artigo sobre
 as matemáticas da física. Algúns pasarán a páxina con 
 preguiza, outros esperarán ecuacións formáis moi bonitas
 e globáis, imposibles de manipular (isto vai polos meus
 pobres estudantes de Métodos II,  compréndovos!). Pois non.
 En absoluto. Este artigo fala de lóxica matemática,
 unha das ramas máis teóricas, descoñecidas, e 
 (desgrazadamente) ignoradas polos demáis profesionáis  
 das matemáticas (salvo excepcións, son coma as bruxas, 
 habelas hainas). É, por así dicilo, a teoría de cordas 
 das mates.

O curioso do asunto é que hai moi boas cabezas pensando 
nas aplicacións desta disciplina á física. Non tanto como 
ferramenta formal, senón para resolver cuestións fundamentáis. 
Un pouco como no caso da teoría de cordas, que xurdiu coa 
motivación de unificar a gravidade e a mecánica cuántica, 
a lóxica matemática, e máis precisamente a subdisciplina 
chamada teoría de modelos, enquírese sobre a base das teorías:  
cal é a linguaxe adecuada?  cales son os axiomas 
(léase postulados) que garantizan o sostén da teoría? 
A moitos sonará o teorema de incompletitude de Gödel, 
segundo o cal, en ``matemáticas", entendida coma a ciencia 
na que os números enteiros son base, sempre haberá enunciados 
indecidibles (é dicir, cuxa veracidade ou falsedade non se 
poda demostrar). 

Pero:  que é a teoría de modelos?  que a fai adecuada para 
esta inmensa tarefa? Pois ben, dunha parte, a teoría de 
modelos ocúpase de calqueira cousa expresable matemáticamente. 
Por exemplo, os grupos abelianos, ou conmutativos, coma 
$\mathbb{Z}$. Para isto, escollemos unha linguaxe que inclúa 
a función +, os cuantificadores $\forall,\, \exists$, e un 
signo especial para o cero. Con isto podemos establecer os 
axiomas ou normas básicas:
\begin{align*}
&\forall x,y\,\, x+y=y+x\\
&\forall y\ 0+y=y+0=y\\
&\forall x\exists y\,\, x+y=y+x=0\\
&\forall x,y,z\,\, (x+y)+z=x+(y+z).
\end{align*}

Dende o punto de vista da teoría de modelos, un grupo abeliano 
é calqueira ``universo'' ou ``estrutura" matemática que 
satisfaga estes axiomas. Dá igual que teña un elemento, 
cinco ou infinitos. Estas simples normas capturan o esencial. 
(Non abandonedes aínda, por favor!). O máis interesante é que 
considerando estes axiomas, o lóxico ten tódolos posibles 
grupos na cabeza ao mesmo tempo. E o mesmo aplica a calqueira 
teoría das matemáticas: para o lóxico non é cuestión deste ou 
dese espazo, senón de todos á vez.
Isto, que pode parecer pouco práctico, da lugar a aplicacións 
impresionantes en análise, por exemplo, ou xeometría. 
E ademáis poderedes contestar aos profesores de métodos cando 
vos critiquen por tachardes o diferencial $dx$. En efecto, na 
teoría dos corpos ordenados, coma os números reáis, existe un 
universo no que o diferencial é un número coma outro calquera, 
máis grande ca $0$, pero máis pequeno que tódolos reáis positivos.  
Así que a tachar diferenciáis sen complexo!

Pois ben, perdidos andaban os lóxicos nas súas disquisicións 
existenciáis, cando Hrushovski e Zilber, dous lóxicos 
brillantes\footnote{O nome do primeiro baraxouse para a Medalla 
Fields, o máis alto galardón en matemáticas (que, por certo, 
ostenta Witten coma único físico na lista). Lamentablemente, 
na miña humilde opinión, Hrushovski non levou este premio. De 
llo ter dado, teríase visibilizado a súa enorme labor e 
achegado a lóxica ao matemático medio.}, deron coas famosas 
estruturas de Zariski. Buscaban \textbf{A} xeometría. \textbf{A 
xeometría}. E deron con algo estrano aos ollos dos comúns, 
é dicir, uns modelos que capturaban á vez a xeometría do 
espazo macroscópico, no que mirar nun orde ou noutro non 
afecta a observación (é dicir, é conmutativa), pero que 
eran non conmutativos en esencia. Coma o mundo cuántico. 
Disto fai moitos anos, a fináis do século XX xa se aplicaban 
estes modelos para probar teoremas importantes, coma a 
conxectura de Mordell--Lang \cite{Ehud_1996}. 

Pero Zilber, apaixoado da física, que perseguía una formulación 
lóxica (en tódolos sentidos) desta ciencia, veuno claro.  
E se podemos aproximar o mundo real, incluindo o mundo 
macroscópico e cuántico,  mediante estas estruturas? 
En colaboración con outros dous lóxicos, Solanki e Sustretov, 
estudaron o caso máis sinxelo: o oscilador armónico 
unidimensional. Hoxe centrareime neste traballo, pero deixádeme 
engadir que Zilber segue no seu empeño de atopar \textbf{A} 
linguaxe (\textbf{A LINGUAXE}) da física dende a teoría de modelos. 
Nos últimos anos, catro artigos prometedores axiomatizan a 
mecánica cuántica (e a unifican coa cuántica estadística) mediante 
a lóxica do continuo \cite{Zilber1_2025,Zilber2_2025,Zilber3_2024,Zilber_2023}. 
Espero comprendelos nalgún momento para explicalo coma hoxe farei co oscilador.

Comecemos polo principio: o oscilador armónico en dimensión 1 (é dicir, 
un peirao cun extremo legado a un punto, o marco, e outro que se move libremente). 
A enerxía total do sistema é 
\[
E=\frac{m}{2}\dot{x}^2+\frac{k}{2}x^2,
\]
onde $V(x)=\frac{k}{2}x^2$ é a enerxía potencial con $k$ a costante de resistencia. 
Polo principio de conservación da enerxía, empregando algún truquiño de integración, 
obténse a ecuación
\[
m\ddot{x}-kx=0,
\]
cuxa solución xeral é $x(t)=\cos(\omega t+\alpha)$, con $\omega=\sqrt{\frac{k}{m}}$. 
Podemos resolver no plano considerando pares $(x(t),\dot{x}(t))$. Se chamamos $p(t)=m\dot{x}(t)$ 
(é dicir, o momento lineal), entón temos
\[
\dot{x}(t)=\frac{p(t)}{m},\qquad\dot{p}(t)=\frac{k}{m}.
\]
Para cuantizar, imos poñer todo en notación Hamiltoniana. Temos as coordenadas canónicas 
$p(t)$ e $q(t)$ (que neste caso son tan só o momento lineal  $p(t)=m\dot{x}(t)$ e a posición 
$q(t)=x(t)$), nas que obtemos as ecuacións de Hamilton (que son as mesmas que antes, pero son 
válidas nun contexto máis amplo, para tódolos sistemas Hamiltonianos)
\[
\dot{p}=\{H,p\}=-\frac{\partial H}{\partial q}
\qquad
\dot{q}=\{H,q\}=\frac{\partial H}{\partial p}
\]
onde $H(p,q)=\frac{p^2}{2m}+\frac{m\omega^2q^2}{2}$ é  o Hamiltoniano, neste caso, a enerxía 
total do sistema,  e a diferencia de signos nas ecuacións débese a xeometría simpléctica subxacente.  
Os corchetes $\{\cdot \}$ chámanse corchetes de Poisson, e teñen sentido para tódalas funcións 
$\{f(p,q),g(p,q)\}$. Por exemplo, se $f=p$ e $g=q$, temos
\begin{equation}\label{eq_corchetes}
\{p,q\}=\frac{\partial p}{\partial q}-\frac{\partial q}{\partial p}.
\end{equation}
Isto é cero, sen dúbida, pero parémonos a mirar un pouco a expresión de \eqref{eq_corchetes}: 
primeiro medimos a variación de $p$ na dirección da posición $q$, e restámoslle a variación de 
$q$ na dirección do momento lineal $p$, e isto da cero; é dicir, medir $q$ e inmediatamente despóis 
$p$, ou medir $p$ e inmediatamente despóis $q$ é o mesmo. Así, no mundo Hamiltoniano clásico, medir 
a posición primeiro e o momento despois é indiferente, son operacións que conmutan.

Á hora de cuantizar, os observables, neste caso $p$ e $q$, pasan a interpretarse como operadores 
actuando no espazo de Hilbert dos estados do sistema, neste caso, o espazo vectorial $\mathbb{C}$, 
adornado cunha estrutura hermítica. A avaliación de operadores interprétase coma a medición. Se 
$P, Q$ son os observables cuantizados, escribimos $[P,Q]=PQ-QP$. Eses corchetes son a cuantización 
dos corchetes de Poisson. Aplicar $PQ$ a un estado do sistema corresponde a observar primeiro a 
posición e, inmediatamente despois, o momento. Sabemos que no mundo cuántico, a observación modifica 
o estado do sistema, más explícitamente,  $[P,Q]=-i\hbar$ onde $\hbar$ é a constante de Planck. Temos 
o Hamiltoniano cuántico
\[H=-\frac{\hbar P^2}{2m}+\frac{m\omega^2Q^2}{2}\]
que pertence á ``álxebra'' de observables (unha forma curta de dicir que combinamos observacións 
mediante consecución e suma). Entón, as solucións $\psi(x)$ da ecuación diferencial 
\[
H\psi(x)=\hbar \omega\left(n+\frac{1}{2}\right) \psi(x)
\] 
para $n\in\mathbb{N}$ comprenden tódolos posibles estados enerxéticos do sistema. Asimesmo, 
podemos pasar dun estado $|n\rangle$ a un estado $|n+1\rangle$ e viceversa mediante os operdores 
de creación $a^\dagger=\sqrt{\frac{m\omega}{2k}}Q-\frac{i}{\sqrt{2\hbar\omega}}P$ e aniquiliación 
$a=\sqrt{\frac{m\omega}{2k}}Q+\frac{i}{\sqrt{2\hbar\omega}}P$ respectivamente. 

Tras este recordatorio, volvemos á lóxica, ainda que de maneira un tanto camuflada. A idea dos tres 
lóxicos é lóxica: para cada $N\in\mathbb{N}$, consideran un espazo $L_N$ consistente nunha recta 
$\mathbb{A}$ (quizáis co infinito $\mathbb{P}^1=\mathbb{A}\cup\{\infty\}$, pensade nos reais co 
infinito, ou os complexos co infinito; en realidade, pensade en todos á vez, porque uns son os 
puntos reáis da recta, e outros os complexos) e, sobre cada punto $x\in\mathbb{A}$,  o grupo de 
raíces $N$-ésimas da unidade (ou $\mathbb{Z}_N$ se o preferides). Estes son os puntos $z$ tal que 
$z^N=1$. Por exemplo, se $N=2$, $z=\pm1$, que observamos que corresponden aos valores $e^{\frac{2k\pi i}{2}}$ 
para $k=0,1$. Para $N$ xeral, obteremos $e^{\frac{2k\pi i}{N}}$ para $k=0,\dots, N-1$. Este grupo 
$\mathbb{Z}_N$ (é un grupo esencialmente porque pódense multiplicar elementos e seguimos dentro, o 
$1$ pertence a $\mathbb{Z}_N$, e ademáis tódolos elementos teñen inverso) representa os posibles 
estadíos do sistema no punto $x$ sobre o que vive, ata enerxía $N$. Para conseguir tódolos estadíos 
enerexéticos haberá que  considerar tódolos $N$'s xuntos; isto non é problemático, pois as construccións 
son compatibles coa relación $N<N+1$, no sentido de que $L_N\subset L_{N+1}$, e as operacións de creación 
e aniquilación de $L_{N+1}$ inducen as de $L_N$. Obsérvese que temos
\[
\pi:L_N=\mathbb{P}^1\times\mathbb{Z}_N\longrightarrow \mathbb{P}^1
\]
que a cada punto $\ell=(x,\gamma)\in L_N$ asígnalle o $x\in \mathbb{P}^1$ sobre o que ``vive''.
Así, a fibra $\pi^{-1}(x)$ representaría os estados do sistema en $x$.

Os operadores de creación (resp aniquilación) pódense modelizar considerando tripletes 
$(\ell,\ell',b)\in A^\dagger\subset L_N\times L_N\times \mathbb{A}$ con $\ell\in\pi^{-1}(x)$, $\ell'\in\pi^{-1}(x+1)$, $b^2=x$ 
(de maneira que $(\gamma\ell',\gamma\ell,b_\gamma)\in A^\dagger$ para todo $\gamma\in\mathbb{Z}_N$ e algún 
$b_\gamma\in\mathbb{A}$)  (resp. $(\gamma\ell',\gamma\ell,b_\gamma)\in A\subset L_N\times L_N\times \mathbb{A}$ con $\ell'\in\pi^{-1}(x+1)$, 
$\ell\in\pi^{-1}(x)$, $b_\gamma^2=x$). Se o triplete $(\gamma\ell,\gamma\ell',b_\gamma)\in A^\dagger$, interpretamos que 
pasamos do estado $\ell$ (que identificamos con algún $|M\rangle$) no punto $x$ ao estado $\ell'\sim|M+1\rangle$ (olliño! 
Tamén no punto $x$). O grupo ven a conto como sigue: en $\log(\gamma)$ pasos, pasaremos do estado $|M\rangle$ ao estado 
$\left|\frac{N}{2\pi i}\log \exp\left(\frac{2\pi iM}{N}+{\log(\gamma)}\right)\right\rangle$. 
A avantaxe desta linguaxe e que non necesitamos identificar $M$, bástanos con saber que subimos 
de estado enerxético en certo número de pasos.

Hai unha certa duplicidade na definición, xa que o feito de crear está na aparición de $\ell$ e $\ell'$ 
no mesmo triplete (sobre puntos diferentes), pero usamos o punto da base para levar a conta de cantos estados 
enerxéticos modificamos. Como basta saber qué estados se modifican nun paso, e despois noutro, etc, basta 
considerar os triples nas fibras $\pi^{-1}(x)\subset L_N$ e $\pi^{-1}(x+1)\subset L_N$. Outra forma de velo 
é que se cortamos (os puntos reáis da) a recta en cachiños $[x,x+1]$, seguimos tendo unha recta dentro. Por 
exemplo:
\[
\mathbb{R}\longrightarrow [x,x+1] \qquad y\mapsto x+1-1/|y-x-1|.
\]
Certo, non está definida en $x+1$, por iso añadimos os infinitos en ambos lados: $x$ identificase co $0$, 
e $x+1$ co infinito. Obsérvese que isto fai perder a referencia de se subimos ou baixamos enerxía, pois 
$[x,x+1]$ e $[x-1,x]$ pasan a ser o mesmo. Para iso está o uso de $b=\sqrt{x}$. Se collemos a raíz positiva, 
subimos, se collemos a negativa, baixamos.

Hai outra construcción posible (e equivalente nun sentido adecuado \cite{Solanki_2013}).  Identificamos, 
dentro de  $L_N\times\mathbb{A}$, os puntos $(\ell\gamma,\gamma^{-1}y)$ para todo $\gamma\in\mathbb{Z}$; 
é dicir, consideramos o conxunto  $[(\ell,y)]=\{(\ell\gamma,\gamma^{-1}y)\colon\gamma\in\mathbb{Z}\}$ 
coma un único punto. Entón, o conxunto $\mathcal{H}_x=\{[\ell,y]\colon \pi(\ell)=x\}$ é un espazo vectorial 
para o produto por escalares $t\cdot [\ell,y]=[\ell,ty]$ e suma 
\[
[\ell_1,y_1]+[\ell_2,y_2]\stackrel{\tiny\exists [\ell_1,y_2']=[\ell_2,y_2]}{=}[\ell_1,y_1]+[\ell_1,y_2']=[\ell_1,y_1+y_2].\footnote{Exercicio: comproba que está ben definido!}\]
sobre o que se poden definir os operadores
\[
a^\dagger[(e,y)]=[(e',yz)]\qquad a[(e',y)]=[(e,yz)]
\]
onde $(e,e',z)\in A^\dagger$, ou $(e',e,y)\in A$. Ben, a cuestión é que considerar $A^\dagger$ ou $a^\dagger$, $A$ ou $a$, 
non cambia a lóxica do asunto, pois as teorías de cada un destes espazos (é dicir, tódolos teoremas satisfeitos 
por cada un deles) son equivalentes nun sentido adecuado que permete obter unha da outra.  

E agora me diredes: pero  isto non é xeometría? onde está a lóxica? Pois a lóxica está en que nada disto se 
pode enunciar formalmente no mundo da xeometría \textbf{complexa} clásica. Pódese se utilizamos os reáis (a clave 
é o uso dunha raíz cuadrada na definición das relacións, cuxas dúas solucións reáis, positiva ou negativa, son 
identificables coa creación ou aniquilación). Pode que isto non pareza revolucionario, pero espérase que os 
exemplos en teoría de modelos veñan da xeometría complexa, e este non o cumple. O que sí que é máis interesante 
é que $L_N$ vive sobre a recta, que \textbf{SÍ} que é complexa, no sentido de que podemos coller puntos reáis ou 
complexos sen cambiarmos nada. Un recubrimento finito dunha recta que non é complexo (coma o $L_N$ aquí presente) 
é \textbf{MOI} insólito. A súa existencia débese á subxacencia de xeometría non conmutativa. O oscilador armónico é un 
tanto tramposo, porque cos operadores de creación e aniquilación (que conmutan) podemos esconder a non conmutatividade 
da posición e momento cuánticos. Pero se miramos a obxectos máis refinados, coma os automorfismos do recubrimento 
$L_N\longrightarrow \mathbb{A}$, veremos que o comportamento é estrano, e codifica a relación $[P,Q]=-i\hbar$. Para 
construirmos espazos desta índole, temos que movernos ao mundo das xeometrías de Zariski puras, fundamentalmente 
modelo-teóricas, 
non provenintes de ningún espazo clásico.  

Recapitulando: a lóxica permítenos construir un espazo, que pode parecer xeométrico, pero non o é, onde vive o 
oscilador armónico cuántico. Aínda que é un exemplo sinxelo, é tamén o primeiro exemplo físico que se consegue 
modelizar completamente con ferramentas lóxicas. Asemesmo, provee un dos primeiros exemplos dunha xeometría de 
Zariski no proveninte da xeometría clásica, o cal engade ao interés destas teorías. Non sei a vós, pero a min 
éncheme de expectación...

\printbibliography
\end{multicols}
